\section{Case study: \texorpdfstring{$\mathbf P^1$}{P1} and every twist \texorpdfstring{$\cO(d)$}{O(d)}}

\subsection{The two-chart atlas}

Fix a commutative ring $R$ with a certified presentation.  The projective line has two standard
charts
\[
   U_0=\Spec R[t],\qquad U_1=\Spec R[u],
\]
with overlap
\[
   U_{01}=\Spec R[t,t^{-1}]=\Spec R[u,u^{-1}],
   \qquad u=t^{-1}.
\]
The overlap maps are localizations at $t$ and $u$, followed by the certified isomorphism sending
$u$ to $t^{-1}$.  The inverse sends $t$ to $u^{-1}$.  Their composite identities are literal
Laurent-monomial identities.

The line bundle $\cO(d)$ is trivial on each chart.  Let $e_0$ and $e_1$ be the chosen local bases.
On the overlap,
\[
                         e_1=t^d e_0.
\]
This formula is valid for every $d\in\Z$ because $t$ is invertible on the overlap.

A degree-zero \v{C}ech cochain is a pair
\[
                   a(t)e_0\in R[t]e_0,
              \qquad b(u)e_1\in R[u]e_1.
\]
After writing $u=t^{-1}$, the differential is
\begin{equation}\label{eq:p1-delta}
  \delta_d(a,b)=a(t)-t^d b(t^{-1})
        \in R[t,t^{-1}]e_0.
\end{equation}
Since the cover has two members, the ordered \v{C}ech complex is
\[
  0\longrightarrow R[t]\oplus R[u]
    \xrightarrow{\delta_d} R[t,t^{-1}]
    \longrightarrow0.
\]

\subsection{Global sections}

\begin{theorem}[Global sections of $\cO(d)$ on $\mathbf P^1$]\label{thm:p1-h0}
The kernel of $\delta_d$ has the following explicit normal form:
\[
 \check H^0(\mathbf P^1_R,\cO(d))\cong
 \begin{cases}
   R^{d+1},& d\geq0,\\
   0,&d<0.
 \end{cases}
\]
For $d\geq0$, the vector $(c_0,\ldots,c_d)$ corresponds to
\begin{align*}
 a(t)&=\sum_{e=0}^d c_e t^e,\\
 b(u)&=\sum_{j=0}^d c_{d-j}u^j.
\end{align*}
The conversion in both directions is primitive recursive on finite polynomial codes.
\end{theorem}

\begin{proof}
The equality $\delta_d(a,b)=0$ says
\[
                         a(t)=t^d b(t^{-1}).
\]
The left side has Laurent support in $\{e:e\geq0\}$; the right side has support in
$\{e:e\leq d\}$.  Equality of Laurent polynomials is coefficientwise, so the common support lies in
\[
                         \{e:0\leq e\leq d\}.
\]
If $d<0$, this set is empty and both polynomials are zero.  If $d\geq0$, the common Laurent
polynomial has unique coefficients $c_0,\ldots,c_d$.  Rewriting its coefficient at exponent $e$ as
the coefficient of $u^{d-e}$ gives the displayed formula for $b$.

No cancellation in $R$ is used; the conclusion follows from the free Laurent-monomial
presentation.  Extraction and reconstruction scan the finite support interval.
\end{proof}

\begin{corollary}[Global regular functions]\label{cor:p1-global-functions}
For every certified commutative ring $R$,
\[
                         \Gamma(\mathbf P^1_R,\cO)\cong R
\]
in the finite-atlas semantics.  The isomorphism sends a constant $c$ to the pair $(c,c)$.
\end{corollary}

This is already a genuinely scheme-theoretic local-to-global theorem: it is not a statement about a
bounded set of rational points, but about compatible regular functions on two affine charts.

\subsection{First cohomology: the middle-support normal form}

For a Laurent polynomial
\[
                         c(t)=\sum_{e=L}^U c_e t^e,
\]
define three finite pieces:
\begin{align*}
 c_{\geq0}(t)&=\sum_{e\geq0}c_e t^e,\\
 c_{\leq d}(t)&=\sum_{e\leq d}c_e t^e,\\
 c_{(d,0)}(t)&=\sum_{d<e<0}c_e t^e.
\end{align*}
Set
\[
 a(t)=c_{\geq0}(t),
 \qquad
 b(u)=-\sum_{e\leq d}c_e u^{d-e}.
\]
Then
\begin{equation}\label{eq:p1-normal-form}
                     c(t)=\delta_d(a,b)+c_{(d,0)}(t).
\end{equation}

\begin{theorem}[Primitive-recursive $H^1$ normal form]\label{thm:p1-h1}
Every degree-one \v{C}ech cochain has the decomposition
\eqref{eq:p1-normal-form}.  The residual term is supported exactly in the finite interval
$d<e<0$, and two such residual terms differ by a coboundary only when they are coefficientwise equal.
Consequently
\[
 \check H^1(\mathbf P^1_R,\cO(d))\cong
 \begin{cases}
   0,&d\geq-1,\\
   R^{-d-1},&d\leq-2,
 \end{cases}
\]
with basis represented by $t^{d+1},\ldots,t^{-1}$ in the second case.
\end{theorem}

\begin{proof}
The high-support part $c_{\geq0}$ is the $a$ contribution to \eqref{eq:p1-delta}.  For an exponent
$e\leq d$, choose the coefficient of $u^{d-e}$ in $b$ to be $-c_e$.  Then the term
$-t^d b(t^{-1})$ contributes $c_e t^e$.  This proves existence of
\eqref{eq:p1-normal-form}.

The image of the $a$ component is supported in exponents $e\geq0$.  The image of the $b$ component is
supported in exponents $e\leq d$.  Therefore a coboundary has zero coefficient at every exponent
$d<e<0$.  A residual supported in that interval can be a coboundary only when all its coefficients
are zero.  The interval is empty exactly when $d\geq-1$, and otherwise contains
$-d-1$ integers.
\end{proof}

The proof gives algorithms for the quotient map, a section selecting the unique normal form, and
witnesses that their composites are identities.  It is therefore stronger than a mere rank count.

\subsection{Examples}

For $d=2$, a global section is
\[
  a(t)=c_0+c_1t+c_2t^2,
  \qquad
  b(u)=c_2+c_1u+c_0u^2.
\]
There is no $H^1$ obstruction because the interval $2<e<0$ is empty.

For $d=-3$, global sections vanish.  Every Laurent cocycle is cohomologous to a unique expression
\[
                       \alpha t^{-2}+\beta t^{-1},
\]
so $\check H^1\cong R^2$.

The executable companion generated and checked 1,200 random Laurent-polynomial decompositions for
$-7\leq d\leq7$.  The tests use exact integer coefficients and verify reconstruction and support
bounds coefficientwise.

\subsection{iPRA implementation}

A Laurent polynomial code stores a lower exponent, an upper exponent, and a finite coefficient
array, or equivalently a sorted sparse list.  Negative exponents are encoded as signed integers.
The normal-form constructor scans the finite support once and writes each coefficient to one of
three output lists according to the decidable comparisons $e\geq0$ and $e\leq d$.

A quantifier-free invariant after processing a support prefix states:
\begin{enumerate}[label=(\roman*)]
\item the input prefix equals the sum of the partial coboundary and partial residual;
\item the chart-zero list contains only nonnegative exponents;
\item the chart-one list contains only nonnegative $u$-exponents;
\item the residual contains only exponents $d<e<0$.
\end{enumerate}
At the end, coefficientwise uniqueness is proved by a bounded scan of the common finite support
interval.  No induction formula quantifies over polynomials or cohomology classes.

\subsection{From \v{C}ech groups to conventional cohomology}

The standard two-chart cover and its intersection are affine, and $\cO(d)$ is quasi-coherent.  The
usual affine-acyclicity comparison identifies the computed \v{C}ech groups with conventional sheaf
cohomology.  The finite arithmetic proof above is independent of that abstract comparison: it
computes the exact complex and its normal forms.  A full internal development of derived-functor
cohomology would add the comparison theorem as a later layer, not alter the finite calculation.
